\subsection{Soundness of sumcheck over a tensor algebra}\label{sec:dpr:sumcheck}
The ring $A$ is not a field, so the
ordinary polynomial root bound over
an integral domain cannot be invoked
without a proof. The following
elementary lemma is the precise
ring-theoretic substitute needed
for Diamond--Posen's sumcheck
(cf.~\cite[Lemma~4.5]{DiamondPosen2024}).

\begin{lemma}[Factor theorem for a commutative ring]\label{lem:dpr:factor}
Let $R$ be a commutative ring with identity,
$P\in R[X]$ and $a\in R$.
There exist $Q\in R[X]$ and a constant
remainder $c\in R$ with
\[
 P(X)=(X-a)Q(X)+c,\qquad c=P(a).
\]
If $P\neq0$ and $\deg P=d\geq1$,
then $\deg Q=d-1$.
\end{lemma}
\begin{proof}
For every exponent $j\geq1$,
\[
 X^j-a^j=(X-a)
 \sum_{i=0}^{j-1}X^{j-1-i}a^i,
\]
as follows by distributing and canceling
successive terms in the finite sum.
Write $P(X)=\sum_{j=0}^dp_jX^j$.
Subtracting $P(a)$ gives
$(X-a)$ multiplied by the sum
$\sum_{j=1}^dp_j\sum_{i=0}^{j-1}
X^{j-1-i}a^i$.
This supplies $Q$ and establishes
$c=P(a)$. If $p_d\neq0$, the coefficient
of $X^{d-1}$ in $Q$ is exactly $p_d$:
only the term $j=d,i=0$ contributes
to that degree. Thus $\deg Q=d-1$.
\end{proof}

\begin{lemma}[Units from the embedded challenge field]\label{lem:dpr:units}
For distinct $a,b\in L$,
the element $\varphi_1(a)-\varphi_1(b)$
is a unit of $A$.
\end{lemma}
\begin{proof}
The difference equals
$\varphi_1(a-b)$, by additivity
of $\varphi_1$. Since $a-b\neq0$,
its inverse $(a-b)^{-1}$ exists in $L$.
Multiplicativity and preservation of
the unit show
$\varphi_1(a-b)\varphi_1((a-b)^{-1})=1$.
Hence the difference is invertible in $A$.
\end{proof}

\begin{theorem}[Root bound inside the embedded field]\label{thm:dpr:ring-roots}
Let $P\in A[X]$ be nonzero with
degree at most $d$, where $d\geq0$.
Then
\[
 \boxed{
 |\{z\in L:P(\varphi_1(z))=0\}|\leq d.}
\]
\end{theorem}
\begin{proof}
Induct on $\deg P$. If $\deg P=0$,
$P$ is a nonzero constant and has
no roots. Suppose $\deg P=D\geq1$.
If it has no root of the form
$\varphi_1(z)$, there is nothing to prove.
Otherwise choose one such root $a$.
Lemma~\ref{lem:dpr:factor} gives
$P(X)=(X-\varphi_1(a))Q(X)$ with
$\deg Q=D-1$.
For any other root $b\neq a$,
\[
 0=P(\varphi_1(b))
   =(\varphi_1(b)-\varphi_1(a))
                     Q(\varphi_1(b)).
\]
The first factor is a unit by
Lemma~\ref{lem:dpr:units}, so
$Q(\varphi_1(b))=0$.
The inductive hypothesis gives at
most $D-1$ such other roots.
Including $a$, the total is
at most $D\leq d$.
No assumption that $A$ is an
integral domain was used.
\end{proof}

\begin{corollary}[Uniform challenge root probability]\label{cor:dpr:ring-prob}
For $P\neq0$ of degree at most $d$
chosen independently of a fresh
uniform $z\in L$,
\[
 \Pr_z[P(\varphi_1(z))=0]\leq d/q.
\]
\end{corollary}
\begin{proof}
Theorem~\ref{thm:dpr:ring-roots}
allows at most $d$ values of $z$
out of the $q$ equally likely
field elements.
\end{proof}

\begin{definition}[True partial-sum polynomial]\label{def:dpr:true-partial}
For each $0\leq i<h$ and each
previously sampled challenge prefix
$z_{<i}=(z_0,\ldots,z_{i-1})\in L^i$,
define
\[
 G_i(X):=
 \sum_{b\in I_{h-i-1}}
 H(\varphi_1(z_0),\ldots,
   \varphi_1(z_{i-1}),X,b)
 \in A[X].
\]
The notation for $i=0$ means
the empty prefix.
\end{definition}

\begin{lemma}[True sumcheck degree and endpoint identity]\label{lem:dpr:true-partial}
For each $i$, the polynomial $G_i$
has degree at most $2$, and
\[
 G_i(0)+G_i(1)
 =\sum_{b\in I_{h-i}}
 H(\varphi_1(z_{<i}),b).
\]
\end{lemma}
\begin{proof}
By Lemma~\ref{lem:dpr:H-degree},
each variable of $H$ has degree
at most $2$. Specializing earlier
variables and later Boolean variables,
then summing, cannot increase the
degree in $X$.
The right-hand finite sum splits
uniquely according to its first bit
$0$ or $1$; its two parts are
$G_i(0)$ and $G_i(1)$, respectively.
\end{proof}

\begin{definition}[Degree-two tensor sumcheck]\label{def:dpr:sumcheck}
Given public initial claim $S_0\in A$,
let $s_0=S_0$. For $i=0,\ldots,h-1$:
the prover sends a polynomial
$g_i(X)\in A[X]$ of degree at most $2$
\emph{before} the verifier samples
$z_i\in L$ uniformly. The verifier
requires
\[
 s_i=g_i(0)+g_i(1),
\]
then sets
$s_{i+1}:=g_i(\varphi_1(z_i))$.
At the end it checks
$s_h=H(\varphi_1(z_0),\ldots,
                  \varphi_1(z_{h-1}))$,
provided the target evaluation can
be obtained from the underlying
large-field evaluation oracle.
\end{definition}

\begin{lemma}[A false running claim forces a nonzero discrepancy]\label{lem:dpr:false-prefix}
Fix a legal round $i$ and a previously
sampled prefix for which
\[
 s_i\neq
 \sum_{b\in I_{h-i}}
 H(\varphi_1(z_{<i}),b).
\]
If the endpoint condition
$s_i=g_i(0)+g_i(1)$ passes,
then the polynomial
$g_i-G_i\in A[X]$ is nonzero
and has degree at most $2$.
\end{lemma}
\begin{proof}
Lemma~\ref{lem:dpr:true-partial}
shows that the Boolean endpoint sum
of $G_i$ is the true running sum.
If $g_i=G_i$, the accepted endpoint
condition would imply that the
claimed $s_i$ equals this true sum,
contrary to hypothesis. Thus
$g_i-G_i\neq0$.
Both polynomials have degree
at most $2$ by the protocol and
Lemma~\ref{lem:dpr:true-partial};
their difference has degree
at most $2$ or is zero, and we
have excluded the latter.
\end{proof}

\begin{lemma}[One-round false-to-true transition bound]\label{lem:dpr:transition}
Under the hypotheses of
Lemma~\ref{lem:dpr:false-prefix},
conditioned on the fixed prefix and
the prover message $g_i$,
\[
 \Pr_{z_i}\left[
 s_{i+1}=
 \sum_{b\in I_{h-i-1}}
 H(\varphi_1(z_{\leq i}),b)
 \right]\leq 2/q.
\]
\end{lemma}
\begin{proof}
The true next running claim is
$G_i(\varphi_1(z_i))$, by
Definition~\ref{def:dpr:true-partial}.
The new claimed value is
$g_i(\varphi_1(z_i))$.
Their equality means that
the nonzero degree-at-most-two
polynomial $g_i-G_i$ vanishes
at $\varphi_1(z_i)$.
The challenge is independent
of the fixed prefix and polynomial,
so Corollary~\ref{cor:dpr:ring-prob}
bounds that probability by $2/q$.
\end{proof}

\begin{theorem}[Tensor sumcheck soundness]\label{thm:dpr:sumcheck-sound}
Suppose $S_0\neq
\sum_{b\in I_h}H(b)$ and the prover's
messages are fixed before their
respective fresh challenges.
Then the probability that all
degree-two sumcheck endpoint conditions
pass \emph{and} the final true
equality $s_h=H(\varphi_1(z))$
holds is at most
\[
 \boxed{\min\{1,2h/q\}.}
\]
This statement holds for arbitrary
adaptive, randomized causal provers.
\end{theorem}
\begin{proof}
The initial running claim is false
by hypothesis and the definition
of $s_0$. If every endpoint condition
passes and the final true equality
holds, there must be a first round
at which the false running claim
becomes true. At every fixed prefix
with a false claim and a passing
endpoint check,
Lemma~\ref{lem:dpr:transition}
bounds the probability of that
transition by $2/q$,
conditional on the already chosen
prover message. A conditional
union bound over the $h$ rounds
gives probability at most $2h/q$.
For a randomized prover, condition
on its private random tape and
average using the finite law of total
probability, as in
Lemma~\ref{lem:security:randomized-prover}.
The result is also at most $1$.
\end{proof}
